% Exact quarantined release candidate for P06.2.
% This file is not a live-book input before independent and production
% verification.  No novelty or priority claim is made.

\begin{proposition}\label{prop:p062-rational-centred-traces}
Let \(K\) be a complete non-Archimedean field with nontrivial absolute
value, let \(n\geq 1\), let
\(\mathbf R=(R_1,\ldots,R_n)\in |K^\times|^n\), and put
\[
 \begin{aligned}
 D^-_{\mathbf R}
  &=\{x\in(\mathbf A^n_K)^{\an}:|T_i(x)|<R_i
                       \text{ for every }i\},\\
 D^+_{\mathbf R}
  &=\{x\in(\mathbf A^n_K)^{\an}:|T_i(x)|\leq R_i
                       \text{ for every }i\}.
 \end{aligned}
\]
For \(a\in K^n\) and a positive polyradius
\(\mathbf r=(r_1,\ldots,r_n)\), write
\(\eta_{a,\mathbf r}\) for the corresponding Gauss point.

If \(h\in C^0(D^-_{\mathbf R},\mathbf R)\) satisfies
\begin{equation}\label{eq:p062-arbitrary-field-zero-current}
                       \mathrm d'\mathrm d''[h]=0
\end{equation}
as a Chambert--Loir--Ducros current, then
\begin{equation}\label{eq:p062-rational-centred-radial-identity}
                         h(\eta_{a,\mathbf r})=h(a)
\end{equation}
for every \(a\in D^-_{\mathbf R}(K)\) and every
\(0<r_i<R_i\).  Consequently \(h\) is constant on
\(D^-_{\mathbf R}(K)\), on the union of all \(K\)-centred Gauss
skeleta, and on the closure of that union.

If \(h\in C^0(D^+_{\mathbf R},\mathbf R)\) and
\eqref{eq:p062-arbitrary-field-zero-current} holds on the absolute
analytic interior of \(D^+_{\mathbf R}\), then the same conclusions hold
on \(D^+_{\mathbf R}\); in this case
\eqref{eq:p062-rational-centred-radial-identity} is valid for
\(a\in D^+_{\mathbf R}(K)\) and \(0<r_i\leq R_i\).
\end{proposition}

\begin{proposition}\label{prop:p062-explicit-base-change}
Keep the notation of \cref{prop:p062-rational-centred-traces}.

\begin{enumerate}
\item If there is a complete algebraically closed extension \(L/K\) such
that
\[
             \mathrm d'\mathrm d''[h_L]=0,
             \qquad h_L=h\circ p_L,
\]
on \(D^-_{\mathbf R,L}\), then \(h\) is constant on
\(D^-_{\mathbf R}\).

\item If \(h\in C^0(D^+_{\mathbf R},\mathbf R)\) and the analogous
equation holds on the absolute analytic interior of
\(D^+_{\mathbf R,L}\), then \(h\) is constant on
\(D^+_{\mathbf R}\).

\item The hypothesis of the first item is automatic if both \(h\) and
\(-h\) are plurisubharmonic in the approximation sense of
Chambert--Loir--Ducros on \(D^-_{\mathbf R}\).  Hence every such \(h\) is
constant.
\end{enumerate}
\end{proposition}
